% Pursuing Stacks (cotes 134-2 to 134-8): archive cote 134-8, site page 36 (PDF page 37), typescript page 590.
% typescript. Draft compared: ps7.tex:1679-1733.
% AI-transcribed by Claude Opus 5.5 (claude-opus-5-5, adaptive thinking, effort high), Message Batches API, batch ps4 (msgbatch_012Aoc1NuskATPsyBeKSYK7J), 2026-11-12; not yet checked by eye.
\scanpage{134-8/36}
\mspage{590}

to exist)  that in \hand{$\mathcal{M}$} there should be enough projectives resp. injecti-%
ves !  This gives what is needed\inserted[typed]{surely\hand{,}} in order to grind through a mutually
dual treatment of relative homology Lf$_{!}$  and cohomology Rf$_{*}$ .  I have
the feeling that the little work ahead\hand{\unsure{,}}for defining the basic operati-%
ons and working out the relevant formulaire,\inserted{(including ``cap products''),} is a matter of mere
routine,  and I do'nt really expect any surprise may come up.  There-%
fore I do'nt feel like grinding through this,  and rather will feel
free to use whenever needed the most evident formulae,  as the adjunc-%
tion formulae between the three functors
\[ \mathrm{Lf}_{!}\ ,\ f^{*}\ ,\ \mathrm{Rf}_{*}\ \ , \]
transitivity isomorphisms for a composition of maps,  possibly also
various projection formulae  -  trying however to be careful with
base change questions,  notably for Lf$_{!}$ ,  and not repeat the same
blunders !

In retrospect,  the main role\inserted[typed]{for me\hand{,}} of the reflections of chapter V
on abelianization has been to become a little more familiar with the
\uline{homology} formalism,\inserted[typed]{in (Cat)} namely essentially with the Lf$_{!}$ operation, which
has been more or less a white spot in my former experience,  centered
rather upon cohomology.  \xout{This} An interesting,  still somewhat routine
byproduct of these ponderings has been the careful formulation of the
duality relationship \xout{\illegible{1}} around the pairings between sheaves and
co-sheaves, namely the operations $*_{k}$ and $\hand{\unsure{\circledast}}_{k}$ .  Granting this,  the
game with integrators and cointegrators \xout{\unsure{is just the most}}\inserted[typed]{boils down to the} standard
reflex,  of taking projective resolutions of the "unit" sheaf or
cosheaf,  k$_{A}$ or k$_{A^{o}}$  - which are the most obvious objects at hand of
all,  \xout{\illegible{2}} in our coefficient categories.

The  one idea which seems to me of wider scope and significance,
in this whole reflection on abelianization,  is the "further step in
linearization",  whereby the models in (Cat) are replaced by their
k-additive envelope,  endowed moreover with their natural \uline{diagonal}
\uline{map}\inserted{(section 109),} .  The psychological effect of this discovery has been an immedia-%
te one  -  \xout{\illegible{1}} it triggered at once the reflection on schematization
of chapter VI.  This reflection,  apparently,  took me into a rather
different direction from those\hand{$\vert$}"k-coalgebra structures in (Cat)" ,
which had reawakened and made more acute the feeling that homotopy
types should make a sense "over any ground ring".  \over{\unsure{c}}{C}oming back now to
the homology and cohomology formalism within (Cat),  it remains \unsure{a}

\margin{36}
